\section*{September 21, 2026: Changing characteristic premiums and repeat sales}
\textit{Drafted by Codex; checks requested by Jacob.}

Allowing house-characteristic premiums to vary by year reduces the pooled
dollar estimate from \$34,851 to \$18,541 on the exact same transactions.
The log estimate moves from 3.0\% to $-1.2\%$, while PPML changes little.
The 2008--2018 half-mile single-family sample, school definitions, cleaning,
price deflator, transaction weights, and school-site clustering are unchanged.
Pooled estimates compare 2014--2018 with 2008--2012 and omit 2013.

\begin{center}
\begin{tabular}{lrrr}
 & Dollar OLS & Log OLS & PPML \\
All sales, fixed premiums & \$34,851 & 3.0\% & 2.1\% \\
All sales, premiums by year & \$18,541 & $-1.2\%$ & 1.3\% \\
Repeat sales, site effects & \$12,269 & $-6.0\%$ & $-5.7\%$ \\
Repeat sales, parcel effects & $-\$1{,}505$ & $-9.0\%$ & $-7.5\%$ \\
\end{tabular}
\end{center}

With year-specific premiums, the 95\% dollar interval runs from
$-\$15{,}750$ to \$52,831. The log interval runs from $-12.9\%$ to 12.2\%.
This specification interacts the entire existing characteristic block with
calendar year, retaining the same transformations and site and year effects.
The change demonstrates sensitivity to that restriction, without identifying
which characteristic or neighborhood process explains it. The 2010 log
coefficient remains $-0.210$ relative to 2012, with an interval excluding zero.
The joint log pre-test changes from p=0.031 to p=0.121; this does not establish
parallel trends.

Repeat sales provide a usable but smaller comparison: 132 treated parcels
at 17 sites and 630 control parcels at 38 sites sell in both 2008--2012 and
2014--2018. Their 1,808 transactions enter annual models; omitting 2013 leaves
1,778. Both repeat models use exactly those transactions and fixed
characteristic premiums. Moving from all sales to the repeat subset lowers
the estimate before parcel effects are introduced. Parcel effects then remove
fixed differences between properties. The resulting dollar interval is
$[-\$51{,}068,\$48{,}059]$; the log interval is $[-27.2\%,13.8\%]$.
This comparison does not establish a precise zero.
The annual dollar plot still rebounds late: 2018 is \$98,281 above the 2012
reference, but 2008 is also \$102,432 above it. Both individual intervals
exclude zero. The near-zero pooled estimate compares the full pre-period
with the full post-period; it does not remove this annual pattern.

Excluding parcels with recorded changes in controlled characteristics or
improvement records leaves 99 treated and 534 control parcels. Parcel-effects
estimates remain similar: $-\$4{,}718$, $-7.9\%$ in logs, and $-6.6\%$ in PPML,
all with intervals including zero. Stable records do not prove no renovation,
and this sensitivity's dollar events reject the joint pre-test (p=0.019).
Both resale selection and recorded improvements can be post-treatment
outcomes. These restrictions do not create an unselected housing-stock sample.
The apparent 2013 sales collapse in the repeat subset reflects its selection
rule: 2013 transactions must supplement the required pre/post sales. Full
estimation-sample sales instead rise from 141 to 156 treated and from 712 to
887 control between 2012 and 2013. The packet now shows both populations.

{\small Reproduction: run Make in
\texttt{tasks/audits/home\_hedonic\_dynamics/code}. The saved data contain
36 models, 198 treatment coefficients, support, and source hashes, with a
standard report. Six baseline fits replicate. Independent OLS reproduces
the new pooled dollar/log estimates; PPML convergence and score checks pass.
The six-page packet includes annual plots, counts, and pooled intervals.}
